% !TEX root = ../../hw05.tex
\begin{proof}[Solution]
  \textit{(a)} Since $a|x|\ge0$ for every $x\in\R$, the given bound
  implies
  \[
    f(x)\ge a|x|-b\ge-b.
  \]
  Thus the range $f(\R)$ is bounded below. It is nonempty because it
  contains $f(0)$. By completeness of $\R$, its infimum exists as a
  real number. Write
  \[
    m=\inf f(\R),\qquad -b\le m\le f(0).
  \]

  \textit{(b)} We follow the hint by choosing inputs whose outputs
  approach $m$. By the defining property of the infimum, for every
  $n\in\N$ there is an $x_n\in\R$ such that
  \[
    m\le f(x_n)<m+\frac1n.
  \]
  Indeed, if no such point existed, $m+1/n$ would be a lower bound
  for $f(\R)$ larger than $m$. These inequalities give $f(x_n)\to m$.

  We next show that the inputs $x_n$ stay in a bounded interval.
  Combining the coercive bound with the choice of $x_n$ gives
  \begin{align*}
    a|x_n|-b&\le f(x_n)<m+\frac1n\le m+1,\\
    a|x_n|&<m+b+1,\\
    |x_n|&<\frac{m+b+1}{a}.
  \end{align*}
  Since $m\ge-b$ and $a>0$, the number
  $R=(m+b+1)/a$ is positive and finite. Thus $(x_n)$ is contained
  in $[-R,R]$. The coercive bound prevents inputs with nearly
  minimal outputs from escaping to infinity, as illustrated in
  \autoref{fig:hw05:coercive-minimum}.

  \begin{marginfigure}
    \centering
    % !TEX root = ../../problems/hw05.tex
% An exact example, f(x)=1+x^2 >= |x|-1. Here m=1 and R=3.
% Shared physical dimensions for interval and number-line diagrams.
\tikzset{
  interval diagram/.style={
    x=1cm,y=6mm,font=\footnotesize,
    color=black,line width=0.5pt,line cap=round,line join=round,
    >={Stealth[length=4pt,width=3pt]}
  },
  interval endpoint/.style={
    circle,draw=black,line width=0.5pt,
    minimum size=4pt,inner sep=0pt,outer sep=0pt
  },
  interval open/.style={interval endpoint,fill=white},
  interval closed/.style={interval endpoint,fill=black},
  interval label/.style={inner sep=1pt}
}

\begin{tikzpicture}[interval diagram,x=6mm,y=4.5mm]
  \draw[->] (-3.4,0) -- (3.5,0) node[right=4pt] {$x$};
  \draw[->] (0,-1.5) -- (0,8.1) node[above=4pt] {$y$};
  \node[below left=5pt] at (0,0) {$0$};

  \draw[gray,densely dashed] (-3,0) -- (-3,2);
  \draw[gray,densely dashed] (3,0) -- (3,2);
  \draw[gray,densely dotted] (-3.15,2) -- (3.15,2);
  \node[above=4pt] at (2.25,2) {$m+1$};
  \node[below=6pt] at (-3,0) {$-R$};
  \node[below=6pt] at (3,0) {$R$};

  \draw[<->,gray] (-3.3,2.3) -- (0,-1) -- (3.3,2.3);
  \draw[<->,domain=-2.55:2.55,samples=100]
    plot (\x,{1+\x*\x});
  \foreach \position in {0.85,0.55,0.25}{
    \fill (\position,{1+\position*\position}) circle[radius=1.6pt];
  }
  \node[interval closed] at (0,1) {};
  \node[left=6pt] at (0,1) {$m$};

  \draw (-2.8,-3.6) -- (-1.9,-3.6);
  \node[anchor=west] at (-1.5,-3.6) {$f(x)=1+x^2$};
  \draw[gray] (-2.8,-5.3) -- (-1.9,-5.3);
  \node[anchor=west] at (-1.5,-5.3) {$|x|-1$};
\end{tikzpicture}

    \caption[A coercive bound traps a minimising sequence.]{%
      An example with $f(x)=1+x^2$ and the lower bound
      $|x|-1$. Here $m=1$. If $f(x)<m+1=2$, the lower bound forces
      $|x|<3$. The indicated points approach the minimum inside
      this bounded interval.}
    \label{fig:hw05:coercive-minimum}
  \end{marginfigure}

  By the Bolzano--Weierstrass theorem, there is a subsequence
  $(x_{n_k})$ converging to some $x_0\in\R$. Continuity of $f$ at
  $x_0$ gives
  \[
    f(x_{n_k})\longrightarrow f(x_0).
  \]
  But the same subsequence satisfies
  \[
    m\le f(x_{n_k})<m+\frac1{n_k},
  \]
  so $f(x_{n_k})\to m$ as $n_k\to\infty$. By uniqueness of limits,
  $f(x_0)=m=\inf f(\R)$. Hence $f$ attains its global minimum.
\end{proof}
